\subsection{非隣接2項型}\label{節：非隣接2項間漸化式型}

  非隣接2項間\bookterm{recurrence}{漸化式}とは,例えば
  \[
    a_{n+2}=f(a_n)\qquad(n\ge1)
  \]
  のように,隣り合わない2項の間に関係がある\bookterm{recurrence}{漸化式}である.
  この形では,\,$a_1$から$a_3,a_5,\ldots$が決まり,\,$a_2$から$a_4,a_6,\ldots$が決まるため,奇数番目と偶数番目に分けて考えるとよい.
  部分列の番号を$k$とし,$u_k=a_{2k-1}$, $v_k=a_{2k}$と置くと,対応は次のようになる.
  \[
    \begin{array}{c|cccc}
      k&1&2&3&\cdots\\ \hline
      u_k&a_1&a_3&a_5&\cdots\\
      v_k&a_2&a_4&a_6&\cdots
    \end{array}
  \]
  元の番号$n$が$2$進むと,部分列の番号$k$は$1$進む.
  従って,$u_{k+1}=f(u_k)$, $v_{k+1}=f(v_k)$は,それぞれ隣接する2項の式になる.
  最後は奇数の$n$に$k=(n+1)/2$,偶数の$n$に$k=n/2$を対応させれば,すべての項に戻せる.

  また,\,$a_{n+2}=qa_n$では,\,$a_{n+2}=0\cdot a_{n+1}+qa_n$と見て\bookterm{characteristic}{特性方程式}を使うこともできる.

  次の\bookterm{sequence}{数列}の\bookterm{general}{一般項}を求めてみよう.

  \noindent \bookblocklink{example-nonadjacent-1}{problem}{answer}{【例題】}\par\nobreak
  \begin{enumerate}[label=(\arabic*),parsep=3pt,format=\bookitemlink{example-nonadjacent-1}{problem}{answer}]
    \item $\displaystyle a_1=2,\,a_2=3,\,a_{n+2}=a_n^2$
    \item $\displaystyle a_1=1,\,a_2=2,\,a_{n+2}=2a_n+1$
    \item $\displaystyle a_1=1,\,a_2=3,\,a_{n+2}=4a_n$
  \end{enumerate}

  \noindent\bookblocklink{example-nonadjacent-1}{hint}{problem}{【方針】}\par\nobreak
  \begin{enumerate}[label=(\arabic*),parsep=3pt,format=\bookitemlink{example-nonadjacent-1}{hint}{problem}]
    \item 奇数番目と偶数番目を別々の\bookterm{sequence}{数列}とみなし,それぞれの二乗関係を調べます.
    \item $b_n=a_n+1$とおいて定数項を消してから,奇数番目と偶数番目に分けます.
    \item $a_{n+2}=0\cdot a_{n+1}+4a_n$と見て,\bookterm{characteristic}{特性方程式}を解きます.
  \end{enumerate}

  \noindent\bookblocklink{example-nonadjacent-1}{answer}{problem}{【解答】}\par\nobreak
  \begin{enumerate}[label=(\arabic*),parsep=3pt,format=\bookitemlink{example-nonadjacent-1}{answer}{problem}]
    \item $\displaystyle a_{2k-1}=2^{\,2^{k-1}},\quad a_{2k}=3^{\,2^{k-1}}\quad(k\ge1)$
    \item $\displaystyle a_{2k-1}=2^k-1,\quad a_{2k}=3\cdot2^{k-1}-1\quad(k\ge1)$
    \item $\displaystyle a_n=\frac54\,2^{n-1}-\frac14\,(-2)^{n-1}$
  \end{enumerate}

  \noindent\bookblocklink{example-nonadjacent-1}{solution}{problem}{【解法】}\par\nobreak
  \begin{enumerate}[label=(\arabic*),parsep=3pt,format=\bookitemlink{example-nonadjacent-1}{solution}{problem}]
    \item $u_k=a_{2k-1},\,v_k=a_{2k}$とおくと,
          \[
            u_1=2,\quad v_1=3,\qquad
            u_{k+1}=u_k^2,\quad v_{k+1}=v_k^2
          \]
          となる.したがって,それぞれ順に2乗していくので,
          \[
            u_k=2^{\,2^{k-1}},\qquad v_k=3^{\,2^{k-1}}
          \]
          である.$u_k=a_{2k-1}$, $v_k=a_{2k}$を戻せば,求める二つの式を得る.
    \item $b_n=a_n+1$とおくと,
          \[
            b_{n+2}=a_{n+2}+1=2a_n+2=2b_n
          \]
          であり,$b_1=2,\,b_2=3$である.
          よって,奇数番目と偶数番目で分けると,
          \[
            b_{2k-1}=2^k,\qquad b_{2k}=3\cdot2^{k-1}
          \]
          となる.$a_n=b_n-1$に戻せば,
          \[
            a_{2k-1}=2^k-1,\qquad a_{2k}=3\cdot2^{k-1}-1
          \]
          を得る.
    \item \bookterm{characteristic}{特性方程式}は
          \[
            r^2-4=0
          \]
          であり,解は$r=2,-2$である.したがって,
          \[
            a_n=A\,2^{n-1}+B\,(-2)^{n-1}
          \]
          とおける.$a_1=1,\,a_2=3$を代入すると,
          \[
            A+B=1,\qquad 2A-2B=3
          \]
          より$A=\dfrac54,\,B=-\dfrac14$である.ゆえに,
          \[
            a_n=\frac54\,2^{n-1}-\frac14\,(-2)^{n-1}
          \]
          となる.
  \end{enumerate}

  \subsubsection{一周期分をまとめる}\label{節：周期的な係数}
  先ほどは二つ先の項が直接与えられていた.
  今度は一つ先の項への係数が,奇数回と偶数回で変わる場合を考える.

  \bookstep{例題：二回分をまとめる}
  \[
    a_1=1,\qquad
    a_{n+1}=\begin{cases}
      2a_n+1&(n\text{が奇数}),\\
      3a_n+2&(n\text{が偶数})
    \end{cases}
  \]
  奇数番目を$x_k=a_{2k-1}$と置くと,$a_{2k}=2x_k+1$であり,
  次の更新まで続けて
  \[
    x_{k+1}=3(2x_k+1)+2=6x_k+5
  \]
  を得る.従って$x_1=1$から$x_k=2\cdot6^{k-1}-1$であり,
  \BookMathLayout{\[
    a_{2k-1}=2\cdot6^{k-1}-1,\qquad
    a_{2k}=4\cdot6^{k-1}-1\quad(k\ge1)
  \]}{\[
    \begin{aligned}
      a_{2k-1}&=2\cdot6^{k-1}-1,\\
      a_{2k}&=4\cdot6^{k-1}-1\quad(k\ge1)
    \end{aligned}
  \]}
  となる.偶数番目は同じ二回の更新の途中の値から求められた.
  最初に加えた$1$も次の更新で$3$倍されるため,二回分の付加項は$3\cdot1+2=5$となった.

  \bookstep{一周期への拡張}
  一般に,$a_{n+1}=p_na_n+q_n$で係数が周期$L$をもつなら,
  一周期分の更新を合成する.
  この一般化では部分列を$0$から数え,$x_k=a_{1+kL}$ $(k\ge0)$と置く.
  すると$x_0=a_1$であり,
  \[
    x_{k+1}=Px_k+Q
  \]
  となり,$P$は一周期の倍率の積,$Q$は各付加項がその後の倍率を受けた寄与の合計である.\footnote{最初の一周期の係数を$p_1,\ldots,p_L$, $q_1,\ldots,q_L$とすると,
  $P=\prod_{j=1}^L p_j$, $Q=\sum_{j=1}^L q_j\prod_{i=j+1}^L p_i$である.空積は$1$とする.
  この表示は倍率に$0$を含む場合も成り立つ.}
  途中の項は元の一回ずつの更新から戻せる.
  $P=0$なら$x_k=Q$ $(k\ge1)$,$P=1$なら$x_k=x_0+kQ$である.
  それ以外は,\sectionref*{節：特性方程式型}と同じく不動点からの差を取る.

  \begin{bookpoint}
    \textbf{更新の間隔を選ぶ}\par
    係数が繰り返すなら,一回ごとに追う代わりに一周期分をまとめる.
    係数が周期的でも,項の値が周期的になるとは限らない.
  \end{bookpoint}
  この二回分の計算は,付録の\sectionref*{節：周期係数を行列で見直す}で,
  行列の積からも見直せる.
  符号も変わる例は\bookproblemref{practice-connections}{4}{接続演習問4}へ.
